% ============================================================
% §2.6  Positioning
% ============================================================

\section{Summary and Dissertation Positioning}
\label{sec:synthesis_positioning}
\label{sec:positioning}
\label{sec:open_world_landscape}
\label{sec:related_work_positioning}
\label{sec:perspectives_integration}
\label{sec:ch2_summary}

The chapter's argument can now be stated compactly.
Novelty is a mismatch defined relative to an agent's model.
Agents respond to mismatch either by assimilating it within existing structure or by accommodating it through representational change.
Open-world autonomy requires accommodation whenever the world invalidates the vocabulary, operators, or abstractions the agent depends on.
Many mature methods adapt within fixed representations; many methods that revise representations do so slowly, globally, or under assumptions that limit their deployment.
Foundation models widen the priors available to an agent, but they do not remove the need for structured, localizable representations.
The missing capability, recurring at each level of the argument, is representational extension that is structural, fast, and localized to the point of failure.

This dissertation takes a specific step into that gap.
It does not claim to solve open-world autonomy in full.
Instead, it provides evidence for a narrower position: structured abstractions make novelty more manageable because they determine where mismatch appears, what must be changed, and what can be reused.
A monolithic policy may register novelty only as degraded performance.
A structured agent can sometimes identify the failed operator, executor, object relation, or interaction variable that made the difference.
Structure is what makes accommodation localizable, and localizability is what can make it fast.

The technical chapters develop this position in three ways.
\Cref{ch:benchmarks} introduces environments and metrics for studying structural novelty in sequential decision-making.
\Cref{ch:rapid_learn} studies action abstractions, where symbolic operators and learned executors make it possible to localize an execution failure and turn recovery into a targeted learning problem.
\Cref{ch:generalization} studies interaction abstractions, where representing skills at the level of object interaction rather than robot-specific control allows behavior to transfer across objects, physical parameters, and embodiments.

These abstraction stories are complementary.
Action abstractions help after novelty has caused failure: they expose where the agent's competence no longer matches the world.
Interaction abstractions can prevent some variation from becoming novelty in the first place: a better representation makes the new case fall within existing competence.
Together they support the dissertation's central claim that the structure of an agent's representation is not merely an implementation detail.
It shapes whether mismatch becomes brittle failure, efficient adaptation, or ordinary generalization.

This position is also compatible with approaches that learn or propose new structure.
The dissertation does not argue that all useful abstractions must be hand-designed.
Continual learning, developmental robotics, computational creativity, and foundation models all point toward agents that can build richer representations over time.
The claim here is that such emergence needs an interface into which new structure can be placed, tested, and localized.
Structured abstractions provide that interface.
The next chapter formalizes the setting, the notion of novelty, the agent framework, and the abstraction concepts used by the rest of the dissertation.
